% VI34-DETAILED-PROBLEM-18
\begin{problem}[VI.34.P18: Saturation does not change projective localizations]\label{prob:vi34-18}
Let \(I\subseteq S\) be homogeneous and define
\[
I^{\mathrm{sat}}=I:S_+^\infty.
\]
Prove that \(I^{\mathrm{sat}}\) is homogeneous and that for every homogeneous \(f\in S_+\),
\[
(I^{\mathrm{sat}})_f=I_f.
\]

\textbf{Provenance.} Canonical VI/34 extension.
\end{problem}

\ifdefined\FullProblemDossiers
\begin{solution}
For every \(m\ge0\),
\[
I:S_+^m
=
\{s\in S:sS_+^m\subseteq I\}.
\]
Because \(I\) and \(S_+^m\) are homogeneous, each colon ideal is homogeneous: if
\[
s=\sum_ns_n
\]
and \(sS_+^m\subseteq I\), multiply by a homogeneous element of \(S_+^m\); the component criterion for \(I\)
forces every \(s_n\) to satisfy the same containment. The ideals \(I:S_+^m\) form an ascending chain, and a union
of ascending homogeneous ideals is homogeneous. Thus \(I^{\mathrm{sat}}\) is homogeneous.

The inclusion \(I_f\subseteq(I^{\mathrm{sat}})_f\) is immediate. Conversely, if \(h\in I^{\mathrm{sat}}\), then for
some \(m\),
\[
hS_+^m\subseteq I.
\]
Since \(f^m\in S_+^m\),
\[
hf^m\in I.
\]
After localizing at \(f\), this says \(h/1\in I_f\). Therefore
\[
\boxed{
(I^{\mathrm{sat}})_f=I_f.
}
\]
Saturation changes only information supported in the irrelevant direction that projective charts remove.
\end{solution}
\fi
